% page 321 du fac-simile (image p-320 du scan)
% bloc 162.6 x 250.6 mm ; marge gauche 39.4 mm
\pgc{17.780mm}{0.000mm}{
\l{}
\l{}
\l{}
\l{}
\l{\cel{51}313}
\l{}
\l{\sou{kremacar}. (\sou{trans.)}\cel{2}Reduktar a cindri (objekto o la korpo di}
\l{\cel{3}homo mortinta). - EFIS.}
\l{}
\l{\sou{kremono}. (\sou{tekn.}) Ferajo uzata por klozar od apertar fenestro;}
\l{\cel{3}stango de fero rekta, muntita an un ek la montanti, quan}
\l{\cel{3}onu adsuprenirigas od adinfreirigas per rotacigar mancho}
\l{\cel{3}tale ke onu enirigas singla ek la extremaji aden seruro-}
\l{\cel{3}boko, od ekirigas li de olu. - F.}
\l{}
\l{\sou{krenelo}. Aperturo facita ye intervali en la parapeto di rem-}
\l{\cel{3}paro, di turmo, por lansar projektili sur la enemiko. - EF.}
\l{}
\l{\sou{kreolo}. Persono di raso blanka, naskinta en la kolonii (olim)}
\l{\cel{3}Hispana di Amerika. -\cel{2}Persono qua naskis en la kolonii Euro-}
\l{\cel{3}pano inter-tropike. - DEFIRS.}
\l{}
\l{\sou{kreozoto}. (\sou{kemio)}. Liquido oleatra, korodiva ed antisepsiiva,}
\l{\cel{3}quan onu obtenas per la distilo di gudro lignala. - DEFIRS.}
\l{}
\l{\sou{krepo}. Stofo ek lano delikata, plu o min diafana, kun filo}
\l{\cel{3}ondoza, quan onu preparas per kelke tordar la filo longes-}
\l{\cel{3}ale (warpe), pose plunjante la texuro aden aquo e frotante}
\l{\cel{3}lu per vaxo. - DEFR.}
\l{}
\l{\sou{krepisar}. (\sou{trans.})\cel{2}Indutar (muro) per trulo o balayilo, ye}
\l{\cel{3}strato di gipso, di mortero, quan onu lasas aspera vice}
\l{\cel{3}planigar lu. - F.}
\l{}
\l{\sou{krepitar}. (\sou{netrans}.) Produktar bruiso ek intersequo di mikra}
\l{\cel{3}kraki. - DEFIS.}
\l{}
\l{\sou{krepono.}\cel{2}Stofo frizita quale krepo, ma di qua la texuro esas}
\l{\cel{3}plu dika, por togi advokatala, sutani, e c. - DEFIRS.}
\l{}
\l{\sou{krepuskulo}. Lumo nepreciza quan donas la suno-radii kande li}
\l{\cel{3}reflektesas da la strati supra di atmosfero, t.e. kande}
\l{\cel{3}Suno esas ja sub horizonto, e qua deskreskas segun ke la}
\l{\cel{3}astro desaparas. - EFIS.}
\l{}
\l{\sou{kreso}.\cel{1}\sou{(bot.)}\cel{2}Planto perena, qua formacas genero ek la fami-}
\l{\cel{3}lio \textquotedbl{}kruciferi\textquotedbl{}, qua kreskas alonge la rivereti, an la tere-}
\l{\cel{3}ni inundata. - L.\cel{1}\sou{lepidium nasturtium}. DEFIR.}
\l{}
\l{\sou{kreskar}. (\sou{netra}ns.) I. (1)\cel{1}\sou{(aludante ento organizita, vivoza}).}
\l{\cel{3}Divenar plu granda gradoze, til la fino di sua developo}
\l{\cel{3}normala. (2) (\sou{metaf}.)Divenar plu granda, plu multa, sen-}
\l{\cel{3}cese, gradoze. - II. (\sou{aludante vejetanti)}.\cel{2}Developar su}
\l{\cel{3}naturale en ula regioni. - FIS.}
\l{}
\l{krespo. Rondi de pasto tre dina, ek farino dilutita kun lakto}
\l{\cel{3}ed ovi, quan onu varsas, lejera-strate, aden padelo indutita}
\l{\cel{3}ye grasajo, por koquar lu. - F.}
}
